\section{Stability of the boundary under the Planck window}\label{sec:window}

The exclusion region of Fig.~\ref{fig:map} is computed at the central abundance, $\Omega_\chi h^2=0.1200$. We adopt the abundance window $0.1176\le\Omega_\chi h^2\le0.1224$, a $\pm2\%$ interval around the central value that contains the $\pm1\sigma$ measurement uncertainty, and ask whether motion inside it changes the CMS classification of a node. The relic equation is solved at both edges for every admissible node, and the efficiency is tested with dedicated samples at four of them.

\subsection{The window at every node}

At the two edges the number of nodes passing the validity conditions of Eq.~\eqref{eq:cuts} is $470$ at $0.1176$ and $471$ at $0.1224$, against $470$ at the centre, and the three sets share exactly the central $470$. The single extra record at the high edge sits on the width cut and does not enter the map, which is drawn at the central abundance. The root is refined by bisection until the residual in $\Omega_\chi h^2$ falls below $10^{-5}$, and a node is retained only if $|{\Omega_\chi h^2}/{\rm target}-1|<5\times10^{-3}$, a tolerance carried by the validity flag. The window maps each node into a band of $8$--$223$~MeV in dark mass, with a median of $39$~MeV, always with the higher abundance giving the lighter state; its magnitude is $0.07\%$ of $m_\chi$ at the median and $0.21\%$ at the ninetieth percentile. The coupling moves by $0.14\%$ at the median, $0.55\%$ at the ninetieth percentile and $0.67\%$ at most.

A single branch carries every node. Sampling the relic curve on a $400$-point logarithmic grid at each of the three abundances gives exactly one sign change at every node examined this way, and away from the root the curve never approaches the target more closely than $1.1\%$ of it; within the range we searched, we find no second branch. The yield is $\varepsilon\,\sigma_{\rm vis}$ and, at fixed $g_t$, the cross section tracks the invisible branching ratio, so the window moves the limit by the change in $\mathrm{BR_{inv}}$ plus at most half a per cent from phase space and efficiency — the efficiency term being the one checked at the four nodes of Section~\ref{sec:endpoints}: $0.02\%$ at the median and $0.56\%$ at most across the grid.

Of the three nodes nearest the threshold, 227 has a branching ratio constant to five decimals across the window and does not move, staying at $\mu_{95}=1.0095$; 2061 behaves the same way at $1.0105$; 1008, whose branching ratio changes by $0.05\%$, moves from $1.0003$ to $0.9998$ and $1.0008$ and crosses. That crossing is an order of magnitude below the $2$--$4\%$ Monte Carlo uncertainty of its template, so the classification of 1008 is set by sample size rather than by the window. Every other point lies farther from the boundary than the largest shift its own two edges produce.

\subsection{The four representative nodes}\label{sec:endpoints}

The relic solve fixes the parameters; the efficiency, which is measured from generated events, is tested with dedicated samples at both edges for four nodes: 138 ($M_V=50$~GeV, $g_t=0.1514$), the survivor closest to unity at the time the samples were defined, $\mu_{95}=1.077$; 139 ($M_V=50$~GeV, $g_t=0.1645$), the next, at $1.085$, deliberately adjacent since the two are the only points within $10\%$ of unity in that set; 272 ($M_V=80$~GeV, $g_t=0.1788$), the closest to unity in the $80$--$120$~GeV range, $\mu_{95}=0.800$, where the yield is largest and the mass shift smallest; and 468 ($M_V=140$~GeV, $g_t=0.9397$), a deeply excluded control at $0.0365$. The samples use the same generator chain, the same seven requirements and the same likelihood as the central points; 138 and 139 carry $10^5$ events per edge in two independent halves, 272 and 468 $2\times10^4$. Table~\ref{tab:endpoints} gives the results together with the selection efficiency of each sample, the quantity that limits the comparison.

\subsection{What the endpoint comparison measures}

Three effects separate a centre from an edge. The relic window moves $m_\chi$ by at most $0.23\%$ and the visible cross section by at most $0.5\%$. The finite sample carries a binomial efficiency uncertainty, $\sqrt{(1-\varepsilon)/(N_{\rm gen}\varepsilon)}$, of about $2.4\%$ relative at $10^5$ events and $4.3\%$ at $2\times10^4$, that is $0.04$ and $0.11$ percentage points. The likelihood scales as $\mu_{95}\propto1/N_{\rm SR}$ near the boundary, so a $5\%$ yield difference produces a $5$--$10\%$ change in the limit.

The measured efficiency differences from the centre are $0.98\sigma$ (138, low-$\Omega$), $0.95\sigma$ (138, high-$\Omega$), $0.72\sigma$ (139, low), $1.52\sigma$ (139, high), $2.25\sigma$ (272, low), $0.16\sigma$ (272, high) and $0.24\sigma$/$0.17\sigma$ (468), all consistent with statistics apart from node 272 at the low-abundance edge. That $2.7\sigma$ deficit is already present in the leading MET requirement; the three cards differ only in the random seed, the cross sections agree to $0.3\%$, the central sample reproduces its cutflow on re-analysis, and each sample scatters by $1.5$--$5\%$ when split in half. It is a fluctuation of the event sample, and it is why the comparison resolves shifts of about $10\%$ rather than one per cent.

At these four nodes the window moves the parameters by less than $0.25\%$ in mass and $0.5\%$ in cross section. Node 138 remains on the surviving side: its two edges, $1.231\pm0.060$ and $1.229\pm0.062$, are farther from unity than the centre value $1.077$. Node 139 sits within $1.5\sigma$ of unity at the centre, and its two edges differ from it by less than the sample-to-sample efficiency uncertainty, so its classification is unresolved rather than changed. Node 272 remains excluded: its low-abundance edge, $0.923\pm0.086$, is nominally below unity but only $0.9\sigma$ from it, while its high-abundance edge is $0.863\pm0.080$. Node 468 stays far from the boundary: its limit moves from $0.0365$ at the centre to $0.0380$ and $0.0422$ at the two edges, a $4$--$16\%$ change that leaves all three values deep in the excluded region.

\begin{table}[htbp]
\centering
\caption{Expected limit at the two edges of the Planck window, with the selection efficiency of each sample. Centre values come from the frozen $2\times10^4$-event samples, evaluated with the same likelihood. Edges for 138 and 139 use $10^5$ events in two independent halves; those for 272 and 468 use $2\times10^4$. The errors on $\mu_{95}$ are the RMS of 200 Poisson resamplings of the fifteen bin contents, so they include the fluctuation of the total yield as well as of the shape; the systematic efficiency uncertainty is not included and is quoted separately in the tables.}
\label{tab:endpoints}
\begin{tabular}{lcccccc}
\toprule
& \multicolumn{2}{c}{centre} & \multicolumn{2}{c}{low-$\Omega$ edge} & \multicolumn{2}{c}{high-$\Omega$ edge} \\
\cmidrule(lr){2-3}\cmidrule(lr){4-5}\cmidrule(lr){6-7}
Node & $\varepsilon$ [\%] & $\mu_{95}$ & $\varepsilon$ [\%] & $\mu_{95}$ & $\varepsilon$ [\%] & $\mu_{95}$ \\
\midrule
138 & 1.760 & 1.077 & 1.662 & $1.231\pm0.060$ & 1.665 & $1.229\pm0.062$ \\
139 & 1.575 & 1.085 & 1.646 & $1.012\pm0.049$ & 1.727 & $0.967\pm0.046$ \\
272 & 2.655 & 0.800 & 2.305 & $0.923\pm0.086$ & 2.680 & $0.863\pm0.080$ \\
468 & 4.510 & 0.0365 & 4.560 & $0.0380\pm0.0027$ & 4.545 & $0.0422\pm0.0026$ \\
\bottomrule
\end{tabular}
\end{table}

The relic solve covers all $470$ nodes and the efficiency sample covers four. For the $470$ the window moves one node across the boundary, 1008, by $0.05\%$ and far inside its own statistical uncertainty; that statement rests on the branching-ratio change alone. What the four-node sample adds is narrower: it shows that the efficiency measured at the edges does not differ from the centre beyond its statistical scatter, which is what licenses treating $\varepsilon$ as unchanged in the other columns. It does not verify those columns individually. Reaching a $1\%$ resolution on $\mu_{95}$ at the edges would require of order $4\times10^{2}$ times the events used here at every surviving point.
